\begin{abstract}
Engineering agents can fail at different trusted-state boundaries: deciding whether to act, repairing a design under verifier feedback, recovering from corrupted intermediate state, and ranking feasible designs under an explicit policy. \textbf{DiagBench} evaluates these boundaries with four verifier-grounded probes: credible triage (P1), constrained repair/search (P2), post-trap recovery (P3), and policy-conditioned ranking (P4).
The current VEH instantiation contains \textbf{763} manifest-backed tasks derived from a 209-paper extraction audit, 52 cleaned structure-review anchors, and a shared analytical oracle. Across \textbf{12} complete P1--P4 model runs, no model dominates all stages: model_A leads P1, model_B leads P2, model_M leads P3, and model_C leads P4. Log-derived response-control profiles provide a compact diagnostic account of these reversals and are more stage-specific than model recency or broad reasoning labels alone.
Two audits support this diagnostic account. A 36-task P3 intervention improves model-mean recovery from 50.0\% to 63.2\% when raw trajectory history is replaced by verifier-authored state summaries, while reducing cascade from 28.1\% to 9.5\%. A matched isomorphic probe shows that models recognize or complete feasible designs more reliably than they construct them. A hardened circuit audit, using updated P1/P2 tasks while retaining the original P3/P4 banks, further tests the diagnostic scaffold beyond VEH. DiagBench therefore treats engineering-agent readiness as boundary compatibility, not a single endpoint score.
\end{abstract}
